\chapter*{List of Abbreviations}

%\renewcommand*\acffont{\textit}
%\renewcommand*\acsfont{\textbf}
%\renewcommand*\acfsfont{\underline}


\begin{acronym}
    \acro  {ads-b}   [ASD-B]   {Automatic Dependent Surveillance - Broadcast}
    
    \acro  {asde} [ASDE] {Airport Surface Detection Equipment}
    
    \acro  {ads-b}   [ASD-R]   {Automatic Dependent Surveillance - Rebroadcast}    
    
    \acro  {asias} [ASIAS] {Aviation Safety Information Analysis and Sharing}

    \acro  {a-smgc} [A-SMGCS] {Advanced Surface Movement Guidance and Control System}
    
    \acro  {assa} [ASSA]  {Airport Surface Situational Awareness Display}

    \acro  {assc} [ASSC]  {Airport Surface Surveillance Capability}
    
    \acro  {atc} [ATC] {Air Traffic Control}

    \acro  {cdti} [CDTI] {Cockpit Display of Traffic Information}
    
    \acro  {easa} [EASA] { European Union Aviation Safety Agency}
    
    \acro  {faa} [FAA] {Federal Aviation Authority}
    
    \acro  {faroa} [FAROA] {Final Approach and Runway Occupancy Awareness Display}

    \acro  {faros} [FAROS] {Final Approach Runway Occupancy Awareness Signal}
    
    \acro  {icao} [ICAO] {International Civil Aviation Organization}

    \acro  {mlat} [MLAT] {Multilateration}
    
    \acro  {oi} [OI] {Operational Incident}

    \acro  {OSA} [OSA] {Operational Safety Assessment}

    \acro  {OSED} [OSED] { Operational Services and Environment Definition}

    \acro  {OTH} [OTH] {Other}

    \acro  {papi} [PAPI] {Precision Approach Path Indicator}
    
    \acro  {pd} [PD] {Pilot Deviation}

    \acro  {rel} [REL] {Runway Entrance Lights}

    \acro  {rgl} [RGL] {Runway Guard Lights}

    \acro  {ri} [RI] {Runway Incursions}

    \acro  {ril} [RIL] {Runway Intersection Lights}

    \acro  {ris} [RIS] {Runway Incursion Severity}

    \acro  {risc} [RISC] {Runway Incursion Severity Calculator}

    \acro  {roa} [ROA] {Runway Occupancy Awareness}

    \acro  {rtca} [RTCA] {Radio Technical Commission for Aeronautics}

    \acro  {rwsl} [RWSL] {Runway Status Lights}

    \acro  {spr} [SPR] {Safety, Performance, and Interoperability Requirements}    

    \acro  {ssr} [SSR] {Special Service Request} 

    \acro  {surf ia} [SURF IA] {Enhanced Traffic Situational Awareness on the Airport Surface with Indications and Alerts (SURF A - only alerts, SURF IA - both indications and alerts)}
    
    \acro  {tcas} [TCAS] {Traffic alert and Collision Avoidance System}
    
    \acro  {tis-b} [TIS-B] {Traffic Information System - Broadcast}

    \acro  {thl} [THL] {Take Off Hold Lights}
    
    \acro  {twy} [TWY] {Taxiway}

    \acro  {vasis} [VASIS] {Visual Approach Slope Indicator System}

    \acro  {vhf} [VHF] {Very High Frequency}
    
    \acro  {vpd} [VPD] {Vehicle Pedestrian Incident}
    
    \acro  {rwy} [RWY] {Runway}


    


% Possible package options:
%   - footnote :
%       The option footnote makes the full name appear as a footnote.
%   - nohyperlinks
%       If hyperref is loaded, all acronyms will link to their glossary entry. With the option nohyperlinks these links can be suppressed.
%   - printonlyused
%       Only list used acronyms
%   - withpage
%       In printonlyused-mode show the page number where each acronym was first used.
%   - smaller
%       Make the acronym appear smaller.
%   - dua
%       The option 'dua' stands for “don’t use acronyms”. It leads to a redefinition of \ac and \acp, making the full name appear all the time and  suppressing all acronyms but the explicitly requested by \acf or \acfp.
%   - nolist
%       The option nolist stands for “don’t write the list of acronyms”.
 
%% to modify default formating, you can redefine the formatting for the \acf \acs and ac commands
 


%\section{Different cases}
 %     \verb|\ac{iot}|        & second use                  & \ac{iot}      \\
  %      \verb|\acl{iot}|       & force the long version      & \acl{iot}     \\
   %     \verb|\acs{iot}|       & force the short version     & \acs{iot}     \\
    %    \verb|\acf{iot}|       & force the full version      & \acf{iot}     \\
     %   \verb|\aclp{iot}|      & force the plural version    & \aclp{iot}    \\ %\acp, \acsp, \acfp
      %  \verb|\acfi{iot}|      & force the italic version    & \acfi{iot}    \\
       % \verb|\ac*{iso}|       & first use, don't mark used  & \ac*{iso}     \\ % same with all other commands
        %\verb|\ac{iso}|        & second use                  & \ac{iso}      \\
        %\verb|\aclp{iso}|      & force the plural version    & \aclp{iso}    \\
     %   \verb|\acused{ICANN}|  & mark as used, no printing   & \acused{icann}\\
      %  \verb|\ac{ICANN}|      & second use, as if it is 1st & \ac{icann}    \\
       % \verb|\ac{gpio}|       & acronym used, but not listed& \ac{gpio}     \\  
       % \verb|\ac{gpio}|       & second use                  & \ac{gpio}    
  


% will not be listed, as it is not used
%% Define a custom plural form of an acronym   

%    \acrodefplural{iot}[IoTs]{Internets-of-Things}
    
%% Preferably, put all acronym definitions in one file and load it
%   \input{acros.tex}
% https://ramibaddour.com/2017/01/18/latex-working-with-acronyms/


\end{acronym}